\documentclass[10pt,a4paper]{amsart}
\usepackage{amsmath,amssymb,amsthm,graphicx}
% Retained minimal preamble for the Stacks Project chapter more-algebra.tex (More on Algebra).
% The chapter file pinned-original-source.tex carries only \input{preamble} on its first line, so no
% preamble text lives inside the chapter itself.  The Stacks Project uses one single shared preamble for
% every chapter; the snapshot stages it identically as sources/stacks-chow/preamble.tex and
% sources/stacks-etale-cohomology/preamble.tex.  The block below is copied mechanically, line for line,
% from sources/stacks-chow/preamble.tex lines 14-18 (the Xy-pic setup) and 154-217 (the native theorem
% declarations together with the mathematical macros).  Nothing was added, removed or redefined; the
% wrapper class and the amsmath/amssymb/amsthm/graphicx base packages are supplied by the assembler.
% Chapter: more-algebra.tex -- More on Algebra.
% For commutative diagrams we use Xy-pic
\usepackage[all]{xy}

% We use 2cell for 2-commutative diagrams.
\xyoption{2cell}

% Theorem environments.
%
\theoremstyle{plain}
\newtheorem{theorem}[subsection]{Theorem}
\newtheorem{proposition}[subsection]{Proposition}
\newtheorem{lemma}[subsection]{Lemma}

\theoremstyle{definition}
\newtheorem{definition}[subsection]{Definition}
\newtheorem{example}[subsection]{Example}
\newtheorem{exercise}[subsection]{Exercise}
\newtheorem{situation}[subsection]{Situation}

\theoremstyle{remark}
\newtheorem{remark}[subsection]{Remark}
\newtheorem{remarks}[subsection]{Remarks}

\numberwithin{equation}{subsection}

% Macros
%
\def\lim{\mathop{\mathrm{lim}}\nolimits}
\def\colim{\mathop{\mathrm{colim}}\nolimits}
\def\Spec{\mathop{\mathrm{Spec}}}
\def\Hom{\mathop{\mathrm{Hom}}\nolimits}
\def\Ext{\mathop{\mathrm{Ext}}\nolimits}
\def\SheafHom{\mathop{\mathcal{H}\!\mathit{om}}\nolimits}
\def\SheafExt{\mathop{\mathcal{E}\!\mathit{xt}}\nolimits}
\def\Sch{\mathit{Sch}}
\def\Mor{\mathop{\mathrm{Mor}}\nolimits}
\def\Ob{\mathop{\mathrm{Ob}}\nolimits}
\def\Sh{\mathop{\mathit{Sh}}\nolimits}
\def\NL{\mathop{N\!L}\nolimits}
\def\CH{\mathop{\mathrm{CH}}\nolimits}
\def\proetale{{pro\text{-}\acute{e}tale}}
\def\etale{{\acute{e}tale}}
\def\QCoh{\mathit{QCoh}}
\def\Ker{\mathop{\mathrm{Ker}}}
\def\Im{\mathop{\mathrm{Im}}}
\def\Coker{\mathop{\mathrm{Coker}}}
\def\Coim{\mathop{\mathrm{Coim}}}

% Boxtimes
%
\DeclareMathSymbol{\boxtimes}{\mathbin}{AMSa}{"02}

%
% Macros for moduli stacks/spaces
%
\def\QCohstack{\mathcal{QC}\!\mathit{oh}}
\def\Cohstack{\mathcal{C}\!\mathit{oh}}
\def\Spacesstack{\mathcal{S}\!\mathit{paces}}
\def\Quotfunctor{\mathrm{Quot}}
\def\Hilbfunctor{\mathrm{Hilb}}
\def\Curvesstack{\mathcal{C}\!\mathit{urves}}
\def\Polarizedstack{\mathcal{P}\!\mathit{olarized}}
\def\Complexesstack{\mathcal{C}\!\mathit{omplexes}}
% \Pic is the operator that assigns to X its picard group, usage \Pic(X)
% \Picardstack_{X/B} denotes the Picard stack of X over B
% \Picardfunctor_{X/B} denotes the Picard functor of X over B
\def\Pic{\mathop{\mathrm{Pic}}\nolimits}
\def\Picardstack{\mathcal{P}\!\mathit{ic}}
\def\Picardfunctor{\mathrm{Pic}}
\def\Deformationcategory{\mathcal{D}\!\mathit{ef}}
\usepackage{cleveref}
\crefname{originallemma}{original lemma}{original lemmas}
\Crefname{originallemma}{Original lemma}{Original lemmas}
\makeatletter
\newlabel{more-algebra-lemma-join-koszul-regular-sequences}{{0669}{source}{}{sourcepointer-1}{}}
\newlabel{more-algebra-lemma-join-koszul-regular-sequences@cref}{{[originallemma][0][]0669}{[1][0][]source}{}{}{}}
\makeatother
\crefname{originallemma}{original lemma}{original lemmas}
\Crefname{originallemma}{Original lemma}{Original lemmas}
\makeatletter
\newlabel{lemma-join-sequences-koszul-complex}{{0664}{source}{}{sourcepointer-2}{}}
\newlabel{lemma-join-sequences-koszul-complex@cref}{{[originallemma][0][]0664}{[1][0][]source}{}{}{}}
\makeatother
\crefname{originallemma}{original lemma}{original lemmas}
\Crefname{originallemma}{Original lemma}{Original lemmas}
\makeatletter
\newlabel{homology-lemma-double-complex-gives-resolution}{{0133}{source}{}{sourcepointer-3}{}}
\newlabel{homology-lemma-double-complex-gives-resolution@cref}{{[originallemma][0][]0133}{[1][0][]source}{}{}{}}
\makeatother
\title{Juxtaposing a Koszul-regular sequence with a sequence that is Koszul-regular modulo the ideal it generates}
\author{The Stacks Project Authors}
\date{stacks-project:more-algebra.tex:local-snapshot-20261012; GNU FDL 1.2 or later}
\begin{document}
\maketitle
\setcounter{subsection}{1}
\noindent\textit{Editorial note.} All delivered mathematical text is the author's own, copied from the pinned Stacks Project chapter file more-algebra.tex; no statement or proof step was rewritten, reordered or removed.\par\medskip
\noindent\textit{Editorial note.} The chapter file pinned-original-source.tex carries only \textbackslash{}input\{preamble\} on its first line, so no preamble text lives inside the chapter. The retained file sources/stacks-more-algebra/retained-preamble.tex is a mechanical, line-for-line copy of the project's single shared preamble as staged in the snapshot (Xy-pic setup plus theorem declarations and mathematical macros). Nothing was added, removed or redefined, and the excerpt serves only as preamble, never as a span coordinate.\par\medskip
\noindent\textit{Editorial note.} Both counted spans start at the \textbackslash{}begin line of their environment, so the statement and the author's proof are delivered whole; surrounding prose of the section is deliberately outside the counted spans and is carried as standard prerequisites instead.\par\medskip
\noindent\textit{Editorial note.} The citation fields (source\_id, source\_version, authors, source\_url, license.name, source\_title) are identical to those of every other card this lane has produced from this chapter, so the batch selector treats them as one source.\par\medskip
\noindent\textit{Allowed standard prerequisites (editorial).}\begin{itemize}
\item Koszul regularity of a sequence and the identification of the Koszul complex as a finite free resolution of the quotient. assumed in the statement, source lines 7639-7646. Supplies both hypotheses and the goal.
\item The Koszul complex of a juxtaposed sequence is the total complex of the tensor product of the two Koszul complexes. the author cites Lemma lemma-join-sequences-koszul-complex at proof line 7660. Performs the first identification of the chain.
\item A double complex whose rows are resolutions computes the homology of the total complex. the author cites Homology, Lemma homology-lemma-double-complex-gives-resolution at proof lines 7662-7666. Justifies the first quasi-isomorphism in the displayed chain.
\end{itemize}\medskip
\ifdefined\hypertarget\hypertarget{sourcepointer-1}{}\fi
\noindent\textit{Source reference.} Original lemma 0669: focused-seeds/AGP/sources/stacks-snapshot-20261011/tags.txt, source lines 7928--7928. The lemma's own canonical tag, quoted so that readers can resolve the label inside the delivered text; it points at the counted statement itself and carries no proof content.\par\medskip
\ifdefined\hypertarget\hypertarget{sourcepointer-2}{}\fi
\noindent\textit{Source reference.} Original lemma 0664: focused-seeds/AGP/sources/stacks-snapshot-20261011/tags.txt, source lines 7923--7923. Navigation pointer to a statement the counted proof refers to but does not contain; nothing load-bearing is delegated to it.\par\medskip
\ifdefined\hypertarget\hypertarget{sourcepointer-3}{}\fi
\noindent\textit{Source reference.} Original lemma 0133: focused-seeds/AGP/sources/stacks-snapshot-20261011/tags.txt, source lines 1421--1421. Navigation pointer to a statement the counted proof refers to but does not contain; nothing load-bearing is delegated to it.\par\medskip
\begin{lemma}
\label{lemma-join-koszul-regular-sequences}
Let $A$ be a ring. Let $I$ be an ideal generated by a Koszul-regular
sequence $f_1, \ldots, f_n$ in $A$. Let $g_1, \ldots, g_m \in A$ be
elements whose images $\overline{g}_1, \ldots, \overline{g}_m$ form a
Koszul-regular sequence in $A/I$. Then $f_1, \ldots, f_n, g_1, \ldots, g_m$
is a Koszul-regular sequence in $A$.
\end{lemma}

\begin{proof}
Our assumptions say that $K_\bullet(A, f_1, \ldots, f_n)$ is a finite free
resolution of $A/I$ and
$K_\bullet(A/I, \overline{g}_1, \ldots, \overline{g}_m)$ is a
finite free resolution of $A/(f_i, g_j)$ over $A/I$. Then
\begin{align*}
K_\bullet(A, f_1, \ldots, f_n, g_1, \ldots, g_m)
& = \text{Tot}(K_\bullet(A, f_1, \ldots, f_n) \otimes_A
K_\bullet(A, g_1, \ldots, g_m)) \\
& \cong A/I \otimes_A K_\bullet(A, g_1, \ldots, g_m) \\
& = K_\bullet(A/I, \overline{g}_1, \ldots, \overline{g}_m) \\
& \cong A/(f_i, g_j)
\end{align*}
The first equality by
Lemma \ref{lemma-join-sequences-koszul-complex}.
The first quasi-isomorphism $\cong$ by (the dual of)
Homology, Lemma \ref{homology-lemma-double-complex-gives-resolution}
as the $q$th row of the double complex
$K_\bullet(A, f_1, \ldots, f_n) \otimes_A K_\bullet(A, g_1, \ldots, g_m)$
is a resolution of $A/I \otimes_A K_q(A, g_1, \ldots, g_m)$.
The second equality is clear. The last quasi-isomorphism by assumption.
Hence we win.
\end{proof}

\end{document}
